\documentclass[10pt,a4paper]{amsart}
\usepackage[margin=19mm]{geometry}
\usepackage{amsmath,amssymb,mathtools}
\usepackage[scr=rsfs,cal=euler]{mathalfa}
\usepackage{xfrac}
\usepackage[unicode,hidelinks,bookmarks=false]{hyperref}
\newcommand{\ie}{i.e.\ }
\newcommand{\cf}{cf.\ }
\newcommand{\resp}{respectively\ }
\newcommand{\Subsection}{\subsection}
\providecommand{\nobreakdash}{\nobreak\hbox{-}}
\setlength{\emergencystretch}{2em}
\multlinegap0pt
\usepackage{amssymb}
\usepackage[all]{xy}
\usepackage{tikz}
\usepackage{xcolor}
\usepackage{bbm}
\usepackage{bm}
\usepackage{mathtools}
\newtheorem{theorem}{Theorem}
\newcommand{\fa}{\mathfrak{A}} 
\newcommand{\gs}{\mathfrak{g}^\ast} 
\newcommand{\hh}{\mathcal{H}} 
\newcommand{\pol}{\mathrm{Pol}(\mathfrak{g}^\ast )} 
\newcommand{\qq}{\mathcal{Q}_{\hbar}} 
\title{{\bf Quantization of bounded symplectic domains associated with compact Lie groups}}
\author{Alexey A. Sharapov}
\date{Source: arXiv:2509.05931v2 (CC BY 4.0)}
\begin{document}
\maketitle

\noindent\emph{Source and licence.} {\bf Quantization of bounded symplectic domains associated with compact Lie groups}. Version arXiv:2509.05931v2 (CC BY 4.0), distributed by arXiv under the Creative Commons Attribution 4.0 International licence, http://creativecommons.org/licenses/by/4.0/, as recorded in the arXiv licence metadata for this exact version and shown on the abstract page https://arxiv.org/abs/2509.05931v2. Full licence notice: \texttt{licenses/arxiv-2509.05931-CC-BY-4.0.txt}.\par\smallskip

\noindent\textit{Editorial scope.} Counted unit: the article's theorem theorem (source0.tex lines 1183--1189). The article's complete original proof of it is delivered with it (source0.tex lines 1191--1214, one \texttt{proof} environment of 24 lines). No mathematics is written, repaired or re-derived: the statement and the proof are the article's own lines, in the article's own notation, and no retained line is substituted or re-flowed.  Retained uncounted as the source-local context the proof needs: lines 749--751; lines 794--796.  Nothing was omitted from any retained span, so the package carries no omission record.  No retained line contains a citation, so the wrapper carries no bibliography and no bibliography entry.  No retained line contains a figure environment, an image-inclusion command or drawing code, so no image is delivered and the package declares no asset dependency.  Editorial formatting only, in addition: an AMS \texttt{amsart} preamble that restates the article's own declarations and macros used by the retained lines, namely the environments \texttt{theorem} on separate counters, together with the standard packages those lines need.  The retained environments print in this document as Theorem 1 in order of appearance, whereas the article prints the same items with the numbering of its own full document.  The complete native source of the article as distributed in the arXiv e-print for this exact version is bundled unchanged as \texttt{sources/184612/source0.tex}.
\begin{equation}\label{QR}
\|\qq(a)\psi \|\leq |a|_{R} \cdot\|\psi\|\,,\qquad \forall a\in \fa\,, \quad \forall \psi\in \hh^l\,.
\end{equation}

\begin{theorem}[Schm\"udgen]  \label{T1}
    Let $D=D(d_1,\ldots, d_n)$ be a bounded semialgebraic set. Then $$ S(D)=\langle d_1,\ldots, d_n, \pol^2 \rangle\,.$$ 
\end{theorem}

\begin{theorem}Let $D(d_1,\ldots,d_n)$ be a bounded and quantizable semialgebraic  domain and let $f\in \mathrm{Pol}(\gs)$ be a polynomial strictly positive on the topological closure of 
$D$, i.e., $f|_{\bar D}>0$.  Then, for each $\psi \in \hh^D$ and each $\varepsilon>0$, there exists a constant $c>0$ such  that 
   $$
   \big(\psi, \qq^D(f)\psi\big) >-\varepsilon 
   $$
   for all $\hbar<c$. 
\end{theorem}

\begin{proof}
By Theorem \ref{T1}, any polynomial $f(x)$ positive on $\bar D$ belongs to the semiring $S(D)$, meaning that 
\begin{equation}
    f=\sum_{e_k\in \{0,1\}} f^2_{e_1e_2\cdots e_d}d_1^{e_1}d_2^{e_2}\cdots d_n^{e_n}
\end{equation}
for some polynomials $f_{e_1e_2\cdots e_k}(x)$. Applying the prequantization map, we can bring the operator $\qq(f)$ into the form
$$
\qq(f)=\sum_{e_k\in \{0,1\}} \qq(f_{e_1\cdots e_n})\big(\qq(d_{1})\big)^{e_1}\cdots \big(\qq(d_n)\big)^{e_n}\qq(f_{e_1\cdots e_n})+\hbar \qq(r_\hbar)\,,
$$
where $r_\hbar$ is some polynomial in the $x$'s and $\hbar$. Then, for any normalized state $\psi\in \hh^D$,
\begin{equation}\label{sum}
\begin{array}{c}
  \displaystyle   \big(\psi, \qq^D(f)\psi\big)=\sum_{e_k\in \{0,1\}} \big(\qq(f_{e_1\cdots e_n})\psi, \big(\qq(d_1)\big)^{e_1}\cdots \big(\big(\qq(d_n)\big)^{e_n}\qq(f_{e_1\cdots e_n})\psi\big) \\[5mm]+\hbar \big(\psi, \qq(r_\hbar)\psi\big)\,.
    \end{array}
\end{equation}
By the Schwarz inequality and Rel. (\ref{QR}), 
\begin{equation}
    \big|\big(\psi,\qq(r_\hbar)\psi\big)\big|^2\leq \|\qq(r_\hbar)\|^2\leq \|\qq(r_\hbar)\|^2\leq |r_\hbar|^2_R\,,
\end{equation}
where $R$ is the radius of the domain $D$. Since the operators $\big(\qq(d_1)\big)^{e_1}\cdots \big(\big(\qq(d_n)\big)^{e_n}$ are positive semidefinite on $\hh^D$, the sum over the $e_k$'s in (\ref{sum}) is non-negative, while the last term tends to zero as $\hbar\rightarrow 0$. Therefore, we can take $c =\varepsilon/(b+\varepsilon) \leq 1$, where
$$
b=\max_{0\leq \hbar\leq 1}|r_\hbar|_R\,.
$$
\end{proof}

\end{document}
